\documentclass[11pt,aspectratio=169]{beamer}

\usetheme{Madrid}
\usecolortheme{dolphin}
\setbeamertemplate{navigation symbols}{}
\setbeameroption{hide notes}

\usepackage[utf8]{inputenc}
\usepackage[T1]{fontenc}
\usepackage{graphicx}
\usepackage{booktabs}
\usepackage{amsmath}
\usepackage{siunitx}

\title[Predictive Edge Caching]{Predictive Edge Caching for 5G/6G Networks}
\subtitle{From Classical Locality Theory to Online ML-Based Eviction}
\author[]{Ahmet-Harun \c{C}ak{\i}r \and Beyzanur Y{\i}lmaz \and Hasan Yazgan}
\institute[]{Department of Computer Science and Engineering \\ CSE 476/575 Network Systems Project}
\date{\today}

\begin{document}

\begin{frame}
  \titlepage
  \note{Open with the problem scale: users perceive milliseconds, but infrastructure pays for misses in backhaul and core latency.}
\end{frame}

\begin{frame}{Roadmap}
  \tableofcontents
  \note{Use this as a promise: historical context, system design, hard results, then deployment implications.}
\end{frame}

\section{Why This Problem Matters}
\begin{frame}{A Short History of Caching}
  \begin{itemize}
    \item \textbf{1968 -- Working Set Model (Denning):} memory systems rely on temporal locality.
    \item \textbf{1990s -- Web caching era:} LRU/LFU become practical standards for proxy and CDN layers.
    \item \textbf{5G/6G era:} edge nodes face \textbf{volatile demand}, not stable desktop workloads.
    \item \textbf{Key transition:} from \emph{reactive eviction} to \emph{predictive retention}.
  \end{itemize}
  \vspace{0.3cm}
  \textit{Takeaway: old principles still hold, but non-stationary traffic requires smarter adaptation.}
  \note{Frame this as continuity, not replacement. We stand on classic locality theory, then extend it with forecasting.}
\end{frame}

\begin{frame}{Motivation in One Question}
  \begin{block}{Core Question}
    If cache space is fixed and traffic is changing, \textbf{which object will be requested again soon}?
  \end{block}
  \begin{itemize}
    \item Every miss triggers backhaul/core retrieval and higher end-user latency.
    \item High-value workloads (short video, live clips, social trends) shift by hour and event.
    \item A policy optimized for yesterday can be harmful today.
  \end{itemize}
  \note{Speak in operational terms: each avoidable miss is both a cost and an SLA risk.}
\end{frame}

\section{Limits of Classical Policies}
\begin{frame}{Why LRU/LFU/FIFO Become Fragile}
  \begin{columns}
    \begin{column}{0.55\textwidth}
      \begin{itemize}
        \item \textbf{FIFO:} ignores behavior, only age.
        \item \textbf{LRU:} responds to recent accesses, but cannot distinguish signal vs noise.
        \item \textbf{LFU:} strong on stationary demand; weak under regime shifts due to stale counts.
      \end{itemize}
      \vspace{0.2cm}
      \textbf{Failure mode:} cache pollution during popularity transitions.
    \end{column}
    \begin{column}{0.45\textwidth}
      \begin{figure}
        \centering
        \includegraphics[width=\textwidth]{figures/hit_rate_over_time.png}
      \end{figure}
    \end{column}
  \end{columns}
  \note{Point to the drop zones: those valleys align with distribution changes in the trace.}
\end{frame}

\section{Our System}
\begin{frame}{Predictive Edge Caching Architecture}
  \begin{columns}
    \begin{column}{0.52\textwidth}
      \begin{enumerate}
        \item \textbf{Request Simulator (C++17)}: reproducible synthetic + trace replay traffic.
        \item \textbf{Edge Cache Controller}: byte-aware store + pluggable eviction strategy.
        \item \textbf{ML Predictor (Python/XGBoost)}: online probability forecast of near-future re-requests.
      \end{enumerate}
      \vspace{0.2cm}
      \textit{Design goal: isolate policy logic while keeping replay deterministic.}
    \end{column}
    \begin{column}{0.48\textwidth}
      \centering
      \includegraphics[width=\textwidth]{figures/arch_diagram.png}
    \end{column}
  \end{columns}
  \note{Emphasize engineering value: same trace, same store, only policy changes, so comparison is fair.}
\end{frame}

\begin{frame}{Feature Space: What the Model Sees}
  \small
  \begin{table}
    \centering
    \begin{tabular}{lll}
      \toprule
      \# & Feature & Intuition \\
      \midrule
      1 & \texttt{frequency} & repeated demand in window \\
      2 & \texttt{recency} & how recently object was touched \\
      3--4 & \texttt{iat\_mean}, \texttt{iat\_cv} & regular vs bursty arrivals \\
      5 & \texttt{user\_div} & broad vs narrow audience \\
      6 & \texttt{norm\_size} & byte-efficiency awareness \\
      7 & \texttt{req\_rate} & short-term request intensity \\
      8--9 & \texttt{time\_sin/cos} & cyclic day-time behavior \\
      \bottomrule
    \end{tabular}
  \end{table}
  \note{Tell audience this is intentionally tabular and interpretable; we avoid black-box complexity for this stage.}
\end{frame}

\begin{frame}{Training Strategy and Stability Improvements}
  \begin{itemize}
    \item Sliding window capped at recent 10,000 labeled samples.
    \item Retraining from scratch avoids unbounded tree growth in incremental boosters.
    \item \textbf{Update applied:} retrain only after a full new interval of labels, preventing repeated expensive retrains.
    \item EWMA blending provides robust cold-start behavior before model maturity.
  \end{itemize}
  \begin{block}{Practical outcome}
    Faster, stable online learning with no ``stuck'' behavior in long runs.
  \end{block}
  \note{Mention this was not just theoretical; it addresses a real runtime bottleneck in the project code.}
\end{frame}

\section{Experimental Methodology}
\begin{frame}{Non-Stationary Workload Design}
  \begin{itemize}
    \item 50,000 requests over 10 epochs.
    \item Each epoch activates a disjoint hot pool (\(\approx\)50 files, 80\% traffic concentration).
    \item Intra-epoch popularity follows Zipf(\(\alpha=1\)); inter-arrival rates modulated over time.
    \item Cache budget is fixed, forcing explicit eviction trade-offs.
  \end{itemize}
  \vspace{0.2cm}
  \textit{This setup stresses adaptation, not just steady-state memorization.}
  \note{Call this a ``regime-shift benchmark'' and explain why LFU should be challenged here.}
\end{frame}

\section{Results}
\begin{frame}{Final Outcome: Hit Rate Comparison}
  \begin{columns}
    \begin{column}{0.46\textwidth}
      \begin{table}
        \centering
        \small
        \begin{tabular}{lrrr}
          \toprule
          Policy & Hits & Misses & Hit Rate \\
          \midrule
          \textbf{ML-Driven} & \textbf{28,745} & \textbf{21,255} & \textbf{57.49\%} \\
          LRU & 23,050 & 26,950 & 46.10\% \\
          FIFO & 20,485 & 29,515 & 40.97\% \\
          LFU & 20,431 & 29,569 & 40.86\% \\
          \bottomrule
        \end{tabular}
      \end{table}
      \vspace{0.2cm}
      \textbf{Relative gain vs LRU:} +24.7\% hit-rate improvement.
    \end{column}
    \begin{column}{0.54\textwidth}
      \centering
      \includegraphics[width=\textwidth]{figures/final_comparison.png}
    \end{column}
  \end{columns}
  \note{Pause after stating 57.49\%. Then immediately convert to operational language in next slide: misses avoided.}
\end{frame}

\begin{frame}{Operational Interpretation: Misses Prevented}
  \begin{itemize}
    \item Misses directly map to core-network fetches and added tail latency.
    \item ML-Driven policy prevents:
      \begin{itemize}
        \item 5,695 misses vs LRU
        \item 8,314 misses vs LFU
      \end{itemize}
    \item Benefit is largest at regime boundaries, where classical policies lag.
  \end{itemize}
  \begin{block}{Network implication}
    Better cache intelligence reduces backhaul pressure and improves user QoE simultaneously.
  \end{block}
  \note{Translate this into CFO + operator language: fewer misses means lower transit and better experience.}
\end{frame}

\section{Demonstration}
\begin{frame}{Live Dashboard Improvements for Presentation}
  \begin{itemize}
    \item Increased UI scale, fonts, and chart clarity for projector readability.
    \item Larger KPI cards and thicker plot lines for quick audience interpretation.
    \item More informative run logs (active trace, cache size, history window).
    \item Progress updates tuned for smoother live demo feedback.
  \end{itemize}
  \note{Tell audience the demo is now optimized for visibility in class/room, not only for developer monitors.}
\end{frame}

\section{Conclusion}
\begin{frame}{Conclusion and Forward Path}
  \begin{itemize}
    \item Predictive eviction outperforms reactive baselines on non-stationary traffic.
    \item The gain is not marginal; it materially changes miss volume.
    \item Lightweight tabular ML (XGBoost + online updates) is sufficient to win in this benchmark.
  \end{itemize}
  \vspace{0.2cm}
  \textbf{Next steps:}
  production traces, geographically distributed learning, and cross-cell transfer.
  \note{Close with one sentence: ``Edge caching is no longer only a storage problem; it is now a forecasting problem.''}
\end{frame}

\begin{frame}
  \centering
  \Huge Thank You\\[0.4cm]
  \Large Questions?
  \note{Invite questions on fairness, reproducibility, and deployment constraints first.}
\end{frame}

\section{References}
\begin{frame}[allowframebreaks]{References}
  \footnotesize
  \begin{thebibliography}{99}
    \bibitem{denning} P. J. Denning, ``The Working Set Model for Program Behavior,'' \textit{Communications of the ACM}, 1968.
    \bibitem{breslau} L. Breslau et al., ``Web Caching and Zipf-like Distributions,'' \textit{IEEE INFOCOM}, 1999.
    \bibitem{traverso} S. Traverso et al., ``Temporal Locality in Today's Content Caching,'' \textit{ACM CCR}, 2013.
    \bibitem{xgboost} T. Chen and C. Guestrin, ``XGBoost: A Scalable Tree Boosting System,'' \textit{KDD}, 2016.
    \bibitem{caching5g} T. Li et al., ``Caching as a Service for 5G Networks,'' \textit{IEEE Access}, 2019.
    \bibitem{sadeghi} A. Sadeghi et al., ``Optimal and Scalable Caching for 5G using RL,'' \textit{IEEE JSTSP}, 2018.
    \bibitem{cisco} Cisco, ``Annual Internet Report,'' 2023.
  \end{thebibliography}
\end{frame}

\end{document}